\section{Kiến thức nền tảng}
\label{sec:kien-thuc-nen}

% ============================================================
% MỤC TIÊU CHƯƠNG:
% Cung cấp nền tảng toán học tối thiểu cần thiết để hiểu
% toàn bộ các phương pháp ở Chương 3 và 4.
% Nguồn: Slides Part 1 (Groups, Group Actions, Representations)
% ============================================================

\subsection{Đối xứng và vai trò trong học máy}
\label{subsec:doi-xung}
% Trình bày: Symmetry là gì? Ví dụ trực quan trong thực tế
% (ảnh phân loại, đám mây điểm, đồ thị, neural networks).
% Tại sao encoding symmetry giúp: fewer degrees of freedom,
% data efficiency, OOD generalization.
% Dẫn nguồn: NequIP (3 bậc ít data hơn), scaling laws.

\subsection{Tính Bất biến và Tính Đẳng biến}
\label{subsec:invariance-equivariance}
% Trình bày:
%   - Invariance: f(T(x)) = f(x)
%   - Equivariance: f(T(x)) = T(f(x))
%   - Sơ đồ giao hoán (commutative diagram)
%   - Invariance là trường hợp đặc biệt của Equivariance
%     (scalar output, transformation = identity)
% Ví dụ: shift/rotation invariance trong image classification,
%   permutation equivariance trong node embedding.

\subsection{Lý thuyết Nhóm cơ bản}
\label{subsec:group-theory}
% Trình bày:
%   - Group G: 4 tính chất (associativity, identity, inverse, closure)
%   - Ví dụ: S_n (permutations), SO(3) (rotations), C_n (cyclic/shift),
%     E(3) (Euclidean: rotation + reflection + translation)
%   - Group action: G acts on X as g.x
%   - Group representation: biểu diễn g bằng ma trận ρ(g),
%     yêu cầu ρ(gh) = ρ(g)ρ(h)
%   - Orbit: Orb_G(x) = {g.x | g ∈ G}

\subsection{Biểu diễn nhóm và ý nghĩa với kiến trúc mạng}
\label{subsec:group-representation}
% Trình bày:
%   - Regular vs. irreducible representations
%   - Schur's Lemma (dẫn dắt tới hạn chế của equivariant linear layers
%     ở Chương 3, mục Tensor Methods)
%   - Fourier analysis trên groups (dẫn dắt tới Expanders ở Chương 4)
% Chú ý: không cần quá sâu, mục tiêu là cung cấp vocabulary
%   cho các chương sau.

\clearpage